% ---------------------------------------------------------------------------
% The manuscript body, shared verbatim by the Springer (LNCS) and IEEE
% (IEEEtran) builds. Nothing class-specific may appear in this file.
%
% Each wrapper defines, before \input-ing this:
%   \wfigure / \endwfigure   a figure spanning the full text width
%   \wtable  / \endwtable    the same for tables
%   \tabscale                the font size command tables are set at
%   \para{...}               a run-in bold paragraph heading
%   \citep{...}               mapped to \cite for the LNCS/IEEEtran builds
% ---------------------------------------------------------------------------

\section{Introduction}

A per-cell classifier scores every grid cell of a harmful-algal-bloom
early-warning system independently: the same gradient-boosted tree, applied
row by row, with no mechanism to see an eddy, a front, or the shape of a
bloom spanning several cells. The natural architectural objection is that
this throws away information a spatial model could use, and the natural fix
is a convolutional model that sees the whole field at once. This paper
reports a controlled test of that objection for one operational system and
finds it does not hold, at least not the way it is usually tried: a
single-frame U-Net~\citep{ronneberger2015} given a \NInputChannels{}-channel
ocean-state image (\NChannels{} physical fields, each paired with an
observed-mask channel) loses clearly to the existing per-cell LightGBM
model on \NTestRows{} identical held-out rows, despite the two being scored
through the same metrics on the same test period by construction.

The result matters beyond this one system for two reasons. First, it was
tested rigorously rather than assumed: the existing tabular model was
reloaded and rescored on the spatial model's own exact test rows rather than
cited from a README, closing the gap between "the spatial model scored
lower than a number we remembered" and "the spatial model scored lower than
the same model, run again, on the same water." Second, the negative result
is informative about what to try \emph{next}, not just what not to do: the
architecture plan that scoped this work deliberately ordered a spatial-only
model before a spatial-plus-temporal one (a
ConvLSTM~\citep{shi2015}) specifically so that a negative single-frame
result would argue against the more expensive next step, not merely
postpone it. It does.

\subsection{Contributions}

\begin{enumerate}\itemsep2pt
\item A \textbf{matched, identical-row comparison} between a spatial U-Net
and the shipped per-cell LightGBM model for t+3 harmful-algal-bloom
probability, built so that a fair comparison does not depend on trusting two
numbers measured at different times on different row sets.
\item A \textbf{live re-verification}: rather than reporting the
\OrigDate{} build's numbers as a standing fact, this paper re-ran the full
training-and-comparison pipeline on \AsOf{} and confirms the original
result to three to four decimal places on every metric (\S\ref{sec:results}).
\item A stated \textbf{engineering hazard specific to mixing this stack}:
importing \texttt{torch} and then loading a LightGBM artifact with
\texttt{joblib} in the same process segfaults on this platform, with no
Python traceback --- silent enough that the fix (score the tabular baseline
in a subprocess) had to be designed as part of making the comparison
trustworthy at all, not as an afterthought (\S\ref{sec:methods}).
\item An explicit \textbf{narrowing of the next step}: the result is
evidence against reaching for temporal recurrence next, not merely a
placeholder result to be superseded by one.
\end{enumerate}

\section{Related Work}
\label{sec:related}

\para{U-Net and focal loss} The U-Net encoder-decoder with skip connections
\citep{ronneberger2015} is the standard architecture for dense, per-pixel
prediction from a small amount of training data, which is the regime this
problem sits in: \NTrainDates{} training dates, not the tens of thousands of
images a segmentation benchmark assumes. Focal
loss~\citep{lin2017} was designed for exactly the class-imbalance problem a
\BloomRate{}\% bloom rate creates, down-weighting the loss contribution of
confidently-correct negatives so a handful of hard positive cells is not
swamped; this codebase's own tabular model handles the same imbalance with a
class-weighting scheme, so the two models face an equivalent, not merely
similar, imbalance problem.

\para{Deep learning versus trees on structured environmental data}
Grinsztajn et al.~\citep{grinsztajn2022} show tree-based models remain
state-of-the-art on medium-sized tabular data even accounting for training
cost, and trace the gap to inductive-bias differences rather than raw
capacity. That result is about the same features presented in a different
\emph{representation} (rows versus columns); the comparison in this paper is
a different axis, holding representation architecture fixed to what each
model naturally wants (a dense image for the U-Net, rows for LightGBM) and
asking instead whether \emph{spatial context itself} --- neighbouring cells,
not just more rows of the same per-cell features --- is worth the switch.
Both findings point the same direction for this kind of data without being
the same finding.

\para{ConvLSTM and the ordering of this agenda item} Convolutional
LSTMs~\citep{shi2015} extend a recurrent architecture with convolutional
gates specifically to model spatiotemporal sequences such as radar
nowcasting, and are the natural reach once a single frame is judged
insufficient. This paper's plan ordered the spatial-only U-Net before the
spatial-plus-temporal ConvLSTM for a stated reason, not simplicity alone: if
a full spatial neighbourhood in a single frame does not help, stacking
several such frames in time is unlikely to be the fix, and a negative result
here is therefore evidence against that next step rather than merely a
smaller version of it.

\section{Methods}
\label{sec:methods}

\para{Task and channels} The target is \texttt{bloom\_t3}, the codebase's
one sanctioned forward-shifted label (t+3 harmful-algal-bloom onset),
already fitted train-only by the existing labelling pipeline and reused
unchanged --- the spatial pilot introduces no second forward shift. The
input is a \NInputChannels{}-channel image per date: \NChannels{}
contemporaneous physical fields (chlorophyll, sea-surface temperature,
mixed-layer depth, a chlorophyll gradient, an SST gradient, Okubo--Weiss
vorticity, and a wind-driven upwelling index), chosen to match the
architectural objection directly --- a stratification proxy, frontal-strength
gradients, an eddy-structure term, and the nutrient-supply driver a per-cell
model cannot see the shape of --- each paired with its own observed-mask
channel, since the Arabian Sea grid used here is a rectangle of ocean and
land and a value channel alone cannot distinguish "zero" from "not observed
here."

\para{Grid and architecture} The native grid is \NativeLat{}$\times$\NativeLon{}
cells at 0.25\si{\degree} over the Arabian Sea, zero-padded to
\PadLat{}$\times$\PadLon{} so three stride-2 pooling levels reach a clean
bottleneck without pooling a single-digit dimension; padded cells carry a
permanently-false valid mask and contribute to neither the loss nor any
reported metric. The network is a three-level U-Net with
\BaseChannels{}/32/64/128 channels per level, GroupNorm rather than
BatchNorm (training batches are tens of frames, not hundreds, where
BatchNorm's running statistics are unreliable), and a masked focal loss
($\alpha=\FocalAlpha{}$, $\gamma=\FocalGamma{}$) that excludes padding and
any cell with no label from the gradient entirely rather than down-weighting
it. The trained checkpoint holds \NParams{} parameters.

\para{Matched split and training} Train/validation/test splits use the
identical chronological-fraction cutoffs the tabular model's own final fit
uses, so the test period is the same water at the same dates by
construction: \NTrainDates{} training, \NValidationDates{} validation, and
\NTestDates{} test dates. Training uses Adam (lr~\LearningRate{}, batch size
\BatchSize{}) with early stopping on validation loss (patience
\Patience{} epochs); the reported run stopped after \EpochsRun{} epochs in
\TrainSeconds{}\,s on Apple Silicon MPS --- early stopping fired well
before any wall-clock or hardware limit was approached, though neither
model capacity nor training schedule was swept (\S\ref{sec:limitations}).

\para{Matched-row scoring, and a segfault it was built to avoid} The trained
U-Net's dense per-date predictions are pivoted back into tidy rows at
exactly the coordinates the tabular pipeline scores, using the shared
fusion layer's dense-cube construction in both directions, so the row set
compared is identical rather than merely similar in size. The shipped
LightGBM artifact is reloaded and rescored, and persistence is recomputed,
on that identical row set --- deliberately in a separate process. Measured
directly on this platform: importing \texttt{torch} anywhere earlier in a
process and then loading the LightGBM pipeline with \texttt{joblib}
segfaults (SIGSEGV, no Python traceback), and the reverse import order does
not crash but is unavailable here, since the same process must load
\texttt{torch} to run the U-Net at all. Scoring the tabular baseline in a
subprocess with no \texttt{torch} import is therefore not an optimisation
but the only way the three-way comparison could be trusted to run at all,
and is stated here because a comparison pipeline that silently crashed
midway would have been a second, unrelated way for this result to look
wrong for reasons that had nothing to do with which model forecasts better.

\section{Results}
\label{sec:results}

\begin{wtable}
\centering
\tabscale
\caption{All three models scored on the identical \NTestRows{} test rows.
Higher is better except Brier. Re-run live on \AsOf{}; the \OrigDate{}
measurement this reproduces is quoted in the text.}
\label{tab:comparison}
\begin{tabular}{lrrrr}
\toprule
Model & PR-AUC & ROC-AUC & Brier & TSS \\
\midrule
Spatial U-Net & \UnetPRAUC{} & \UnetROCAUC{} & \UnetBrier{} & \UnetTSS{} \\
LightGBM (reloaded) & \textbf{\LgbmPRAUC{}} & \textbf{\LgbmROCAUC{}} & \textbf{\LgbmBrier{}} & \textbf{\LgbmTSS{}} \\
Persistence & \PersistPRAUC{} & \PersistROCAUC{} & \PersistBrier{} & \PersistTSS{} \\
\bottomrule
\end{tabular}
\end{wtable}

Table~\ref{tab:comparison} was produced by re-running
\texttt{hab\_early\_warning.src.spatial\_train} in full on \AsOf{} ---
training a fresh U-Net from the same seed and configuration, reloading the
same shipped LightGBM artifact, and rescoring both plus persistence on the
\NTestRows{} rows (positive rate \TestPositiveRate{}\%) that run's own test
split produces. Every metric matches the \OrigDate{} account to three or
four decimal places: PR-AUC \UnetPRAUC{} (originally 0.544), ROC-AUC
\UnetROCAUC{} (0.897), Brier \UnetBrier{} (0.058), TSS \UnetTSS{} (0.655)
for the spatial model, and equivalently close agreement for LightGBM and
persistence. This is a second independent training run, ten days after the
first, on a maintained codebase that has continued to change elsewhere in
the interim (\S\ref{sec:limitations}); the reproduction is not a
formality.

\para{Against persistence the picture is mixed; against LightGBM it is not}
The U-Net clears persistence on PR-AUC (gap \UnetVsPersistGap{}) and on
Brier by a wide margin (\UnetBrier{} against \PersistBrier{}), and trails
it marginally on ROC-AUC and TSS (by \(0.004\) and \(0.004\)
respectively) --- a mixed result, not a clean win. The persistence
baseline here is not a bare carried-forward label but a sigmoid-squashed
standardised chlorophyll anomaly, a real, continuous, ranked score; its
comparatively poor Brier alongside a competitive ROC-AUC is consistent with
that score being reasonably well \emph{ordered} but poorly
\emph{calibrated} to the \TestPositiveRate{}\% test-period base rate, not
with it being a degenerate baseline that is trivially easy to beat. Against
the reloaded LightGBM model the comparison is unambiguous: the U-Net trails
on all four metrics simultaneously (PR-AUC gap \UnetVsLgbmGap{}), with no
metric on which it is competitive.

\para{The reloaded baseline agrees with its own published number, which is
what makes the comparison trustworthy} The reloaded LightGBM's PR-AUC,
\LgbmPRAUC{}, sits close to \texttt{hab\_early\_warning/readme.md}'s
documented \ReadmeLGBMPRAUC{} (persistence \ReadmePersistPRAUC{}) for the
same model, though not identical --- that number's row set passed a
different multi-horizon coverage filter, while this run's is t+3-only. The
closeness is itself a check: the comparison pipeline is measuring the
actual shipped model on real held-out water, not something subtly
different that happens to run without error.

\para{The exported artifact is consistent with the region it claims to
cover} A standalone forecast export
(\texttt{exports/hab\_bloom\_forecast\_spatial.nc}, not wired into the
platform's serving path) was re-opened for this paper and its finite-value
fraction measured directly at \ExportFiniteFrac{}\%, matching the
Arabian-Sea grid's independently measured real-ocean-cell fraction of
approximately \ExportOceanFrac{}\% --- land is masked to \texttt{NaN}
correctly, not merely by omission.

\section{Discussion}

Two explanations for a negative result are worth distinguishing, because
they license different next steps. The first is that the implementation is
wrong: a bug, a leak, a mismatched row set. The matched-row, reloaded-
baseline, twice-reproduced design in \S\ref{sec:methods}--\ref{sec:results}
is built specifically to rule this out, and the close agreement with the
tabular model's own published number is the check that it worked. The
second is that spatial context, added this way, genuinely does not carry
enough additional signal for this task to be worth a harder-to-train model
--- and that is the more likely reading here, given the U-Net still clears
persistence by a real margin and is not failing outright, it is simply
behind a strong, already-well-tuned tabular alternative.

\para{What this narrows, not just what it rules out} The architecture plan
this pilot answers to explicitly staged a spatial-only model before a
spatial-plus-temporal one for a stated methodological reason: if a full
neighbourhood in one frame does not help, stacking several such frames
along time is unlikely to rescue the approach, since temporal recurrence
adds a different kind of information (persistence and trend) that a
sufficiently strong tabular model can already partly capture through
lagged features. A negative U-Net result is therefore a specific, informative
data point against reaching for ConvLSTM~\citep{shi2015} next by default,
not an argument that no spatial architecture could ever help this problem.

\section{Limitations}
\label{sec:limitations}

This is one region (Arabian Sea), one horizon (t+3), one channel set, and
one architecture family (a plain U-Net, not a wider search over depth,
channel width, or loss function). The codebase continues to change --- this
paper's own re-verification on \AsOf{} exists precisely because a repository
that trains and retrains models should not be trusted to reproduce a
result from prose alone, a discipline this same research programme applies
to its own other papers~\citep{risingskill2026}. The exported artifact is
explicitly a standalone, unwired output: it is not served to any product
surface, and this paper makes no claim that it should be, absent a future
spatial attempt that actually beats the point model on held-out skill.
Boyce-index spatial calibration, computed for the platform's habitat-
suitability models elsewhere, was not populated for this classification
comparison and is not reported here rather than reported as a placeholder.

\section{Conclusion}

A single-frame U-Net given full spatial context for t+3 harmful-algal-bloom
forecasting beats a naive persistence baseline but loses clearly to the
existing per-cell LightGBM model on \NTestRows{} identical held-out rows,
reproduced live to three to four decimal places ten days after the original
measurement. The result is a real, informative negative: it does not show
spatial structure is irrelevant to blooms, only that this way of giving a
model access to it does not beat a well-tuned tabular alternative on this
task, and it specifically argues against reaching for a spatiotemporal
architecture next, since the ordering of this agenda item was designed to
make exactly that inference available.
